\documentclass{article}
\usepackage{siunitx}
\usepackage{amsmath}
\usepackage{microtype}

\title{Problem Statement}
\date{2026--02--23}
\author{Ahaan Shokeen}

\begin{document}
\maketitle

Redshift and blueshift are phenomena wherein light from distant objects is ``stretched out''; that is, the observed wavelength increases or decreases relative to the emitted wavelength. This affects the light received on Earth, and can be defined by the following equation:
\begin{align}
    z &= \frac{v}{c\sqrt{1 - \frac{v^2}{c^2}}}
\end{align}
wherein \(z\) is the redshift (If \(z = 0\), then there is no redshift or blueshift. If \(z < 0\), it is blueshift. And if \(z > 0\), then there is redshift.), \(v\) is the recessional velocity, in \unit{\metre\per\second}, and \(c\) is the speed of light in a vacuum, approximately \qty{299 792 458}{\metre\per\second}.

This is a problem for cosmologists and spectrographers. Cosmologists need accurate images to study the universe, and whilst tools exist to remedy this, they are nonetheless very slow and inefficient. Spectrographers --- especially extraterrestrial spectrographers --- rely on spectral lines to determine the elemental composition of planets' atmospheres, stars' atmospheres, and more. Our product aims to enhance our interpretation of redshift (Not stop it, because that would require stopping the expansion of the universe, which is beyond our budget.\textsuperscript{[citation needed]}) by accounting for the redshift digitally.

\end{document}
